\section{Introduction}
\label{sec:intro}

Simulation is where most robot learning now happens, and every simulated robot starts as a model
file: an MJCF for MuJoCo~\citep{todorov2012mujoco}, a URDF for ROS-based
stacks~\citep{quigley2009ros}, or a USD for Isaac~\citep{makoviychuk2021isaac}. Writing one is
slow, expert work, and curated collections such as MuJoCo Menagerie~\citep{zakka2022menagerie}
exist because good models are scarce. Large language models have been proposed as robot
designers~\citep{stella2023llmdesign} and used to evolve
morphologies~\citep{ringel2024text2robot,qiu2024robomorph}, but these systems search over
restricted, hand-built parameterisations. Asking a general coding agent to write a new model from
scratch exposes the problem directly: the formats encode every body as a pose relative to its
parent, so a single wrong offset produces a part that floats, a joint that turns about the wrong
axis, or two links that overlap at rest. Nothing in the file reports the error. The agent, which
is otherwise good at fixing what a tool tells it is wrong~\citep{shinn2023reflexion,madaan2023selfrefine,jimenez2024swebench},
has nothing to read.

We take the view that the design interface should make such errors unrepresentable where possible
and loud where not, and should turn every remaining property that matters into a number an agent
can read. \ms implements this view as a Python library with three commitments:

\begin{enumerate}
\item \textbf{Connect, don't place.} There is no API to position a part. A part enters the robot
  only by being fastened face to face to a connector on a part already placed, attached with a
  checked contact, or made the moving side of a joint. Placement is derived, so a robot is
  connected by construction, and a shape that does not touch what it is attached to raises an
  \code{AttachmentError} naming the gap in millimetres (\cref{sec:connect}).
\item \textbf{Everything derived.} Joint anchors, axes, kinematic tree, body frames, inertias and
  actuator limits are computed from the connections, the solids and datasheet specifications; the
  designer never writes them (\cref{sec:derive}).
\item \textbf{A report, not a picture.} A build emits MJCF, URDF and a JSON report of checks: mass
  per body, collisions at home, the usable range of each joint with the first contact that ends
  it, and a test matched to the robot class. Every parameter carries a declared range, so an agent
  can explore without leaving the space the designer validated (\cref{sec:report}).
\end{enumerate}

Our contributions are:
\begin{itemize}
\item A robot description language whose core invariant (every part is reached through a checked
  connection) removes the dominant class of silent errors in hand-written simulation models.
\item A verifier suite that converts physical validity into structured feedback suited to
  agentic iteration, from geometric checks to closed-loop simulation tests.
\item Evidence of expressiveness: six real robots rebuilt from public references, including
  closed kinematic loops, thrust-vectoring rotors and torque-controlled legged platforms, each
  matching its reference's joint axes, and models that train walking policies without further
  editing (\cref{sec:fidelity,sec:rl}).
\item \planned{A controlled study of coding agents building robots with \ms against MJCF/URDF
  authoring and free-placement code CAD, with ablations of each verifier (\cref{sec:agents}).}
\end{itemize}
